\section{Exact admission transport}
\label{sec:terminal-modal-admission}

The optimum transport can be written without a linear solve.  This exposes
the one scalar increment that remains to be controlled in the admission
recurrence.  Let $p_n^\star$ denote the normalized restricted optimum on the
proper prefix $\{v_0,\ldots,v_{n-1}\}$, where $1\le n\le m$, and put
\begin{equation}
 \chi=\frac{1-q}{1+q},\qquad
 f_j=\chi^j+\chi^{-j},\qquad
 b_n=\frac{\rho-q\chi^n}{\chi^n+\chi^{-n}},\qquad
 b_0=\frac{\rho-q}{2}.
 \label{eq:terminal-modal-prefix-parameters}
\end{equation}

\begin{lemma}[Exact prefix optima and transport; proved here]
\label{lem:terminal-modal-prefix-optima}
For $1\le n\le m$ and $0\le j<n$,
\begin{equation}
 p_n^\star(j)=q\chi^j-\rho+b_nf_j.
 \label{eq:terminal-modal-prefix-optimum}
\end{equation}
If the formula on the right is evaluated at the ghost coordinate $j=n$,
its value is zero.  Consequently, for $1\le n\le m$,
\begin{equation}
 p_n^\star-[p_{n-1}^\star;0]
 =(b_n-b_{n-1})f\big|_{\{0,\ldots,n-1\}},
 \label{eq:terminal-modal-prefix-displacement}
\end{equation}
where $p_0^\star$ is empty.  Moreover,
\begin{equation}
 0<b_n-b_{n-1}
 \le q\operatorname{artanh}q
 \le\frac{q^2}{1-q^2}.
 \label{eq:terminal-modal-prefix-increment}
\end{equation}

The full-path optimum, whose far endpoint has ambient degree one, is
\begin{equation}
 p_{m+1}^\star(j)
 =\frac{q(\chi^j+\chi^{2m-j})}{1-\chi^{2m}}-\rho,
 \qquad 0\le j\le m.
 \label{eq:terminal-modal-full-optimum}
\end{equation}
Thus the final displacement from the last proper prefix is also explicit: if
$A_\star=q/(1-\chi^{2m})$ and
$B_\star=q\chi^{2m}/(1-\chi^{2m})$, then
\begin{equation}
 p_{m+1}^\star-[p_m^\star;0]
 =(A_\star-q-b_m)\chi^j+(B_\star-b_m)\chi^{-j},
 \quad 0\le j\le m.
 \label{eq:terminal-modal-final-displacement}
\end{equation}
\end{lemma}

\begin{proof}
Set $y_j=p_n^\star(j)+\rho$.  On every two-sided interior row, the
restricted-optimum equation is homogeneous and has characteristic roots
$\chi$ and $\chi^{-1}$, so $y_j=A\chi^j+B\chi^{-j}$.  The seed row gives
$A-B=q$.  The missing edge at the moving frontier is exactly the ghost
condition $y_n=\rho$.  Hence $B=b_n$, $A=b_n+q$, which proves
\eqref{eq:terminal-modal-prefix-optimum}.  The preceding prefix formula has
ghost value zero at the newly admitted coordinate, so subtraction proves
\eqref{eq:terminal-modal-prefix-displacement}.

For the increment, write $b_n=q\beta(\chi^n)$ with
\[
 \beta(u)=\frac{u(1/5-u)}{1+u^2}.
\]
Because $\chi\ge1-2q$ and $n\le m$,
$\chi^n\ge\chi^m\ge1-2mq=7/8$.  On $[7/8,1]$, $\beta'<0$ and
\[
 -u\beta'(u)=\frac{u(2u+(u^2-1)/5)}{(1+u^2)^2}\le\frac12.
\]
Indeed, after clearing the positive denominator, the last inequality is
equivalent to
\[
 (1+u^2)^2-2u\left(2u+\frac{u^2-1}{5}\right)
 =(1-u^2)\left(1-u^2+\frac{2u}{5}\right)\ge0.
\]
Integration from $\chi^n$ to $\chi^{n-1}$ gives
\[
 0<b_n-b_{n-1}\le\frac q2\log\frac1\chi
 =q\operatorname{artanh}q.
\]
The final inequality follows by integrating
$(1-t^2)^{-1}\le(1-q^2)^{-1}$ for $0\le t\le q$.

On the full path the far-endpoint equation gives
$B\chi^{-m}=A\chi^m$, while the seed equation still gives $A-B=q$.
Solving these two equations proves \eqref{eq:terminal-modal-full-optimum};
subtracting the proper-prefix formula, including its zero ghost value at
$j=m$, proves \eqref{eq:terminal-modal-final-displacement}.
\end{proof}

The frontier value of the prefix optimum has uniform bounds that will also
be used in the admission reduction.

\begin{lemma}[Proper-prefix frontier bounds; proved here]
\label{lem:terminal-modal-frontier-bounds}
Put $P_n=p_n^\star(n-1)$.  For $1\le n\le m$, $P_n$ is nonincreasing in $n$
and
\begin{equation}
 \frac{3488}{1921}q^2\le P_n\le\frac{33}{16}q^2.
 \label{eq:terminal-modal-frontier-bounds}
\end{equation}
More precisely, if $t=\chi^n$, then
\begin{equation}
 \frac{P_n}{q^2}=\frac{2}{1+q}
 \frac{t(1+t/5)/\chi+t-1/5}{1+t^2}.
 \label{eq:terminal-modal-frontier-optimum}
\end{equation}
\end{lemma}

\begin{proof}
Direct substitution in \eqref{eq:terminal-modal-prefix-optimum} proves
\eqref{eq:terminal-modal-frontier-optimum}.  The fraction there is increasing
in $t$: after differentiation its numerator is
$(1+\chi^{-1})(1-t^2+2t/5)>0$.  Because $t=\chi^n$ decreases with $n$, this
also proves monotonicity.

For the lower bound, Bernoulli's inequality and $mq=1/16$ give
$t\ge\chi^m\ge(1-2q)^m\ge7/8$.  Dropping the favorable factor
$\chi^{-1}\ge1$, using that
\[
 \frac{2t+t^2/5-1/5}{1+t^2}
\]
is increasing on $[7/8,1]$, and using $q\le1/16$ yield
\[
 \frac{P_n}{q^2}
 \ge\frac{32}{17}\frac{109}{113}
 =\frac{3488}{1921}.
\]
For the upper bound, evaluate the increasing fraction at $t=1$.  The result
is
\[
 \frac6{5(1-q)}+\frac4{5(1+q)}
 \le\frac{864}{425}<\frac{33}{16},
\]
where the middle expression is increasing on $0\le q\le1/16$.
\end{proof}

There is also an exact local factorization behind the admission analysis.
On a two-sided interior stencil let $\mathcal S_\pm$ be the two unit shifts
and
\begin{equation}
 \mathcal R_\pm=\frac{1-q}{2}(I+\mathcal S_\pm).
 \label{eq:terminal-modal-directional-operators}
\end{equation}
If $A_n$ denotes any projection-inactive fixed-face residual sequence, its
two-sided interior recurrence factors as
\begin{equation}
 A_{n+1}=(\mathcal R_++\mathcal R_-)A_n
 -\mathcal R_+\mathcal R_-A_{n-1}.
 \label{eq:terminal-modal-directional-factorization}
\end{equation}
Indeed, expanding the shifts gives the two coefficients in the fixed-face
recurrence: $\mathcal R_++\mathcal R_-$ is its first-order coefficient and
$\mathcal R_+\mathcal R_-$ is its second-order coefficient.  Thus
$E_{n+1}:=A_{n+1}-\mathcal R_+A_n$ obeys
$E_{n+1}=\mathcal R_-E_n$ away from boundaries.

\begin{lemma}[Proper-prefix one-step chronology; open]
\label{lem:terminal-modal-admission-chronology}
For all sufficiently large $m$, at every proper prefix
$\{v_0,\ldots,v_{n-1}\}$, $1\le n\le m$, the raw proximal point is strictly
positive, the post-step active residual has nonpositive maximum, and the
unique exterior row $v_n$ has normalized residual below $-q^3/5$.  Thus
projection is inactive, the safe correction is zero, and the complete gate
admits exactly the next singleton.
\end{lemma}

Here is an exact reduction of that open statement.  Define the
\emph{linear candidate replay} by starting from the same zero state, using
the raw point
\[
 \widetilde p_n=y_n-r_n(y_n),\qquad
 y_n=\frac{p_n^{\rm in}+qv_n^{\rm in}}{1+q},
\]
at every proper prefix, zero-padding $\widetilde p_n$, and applying the stated
center update and exact optimum transport when the next singleton is
admitted.  This is a recursive algebraic trajectory; it does not assume that
the projected algorithm follows it.  On this candidate put
\begin{equation}
 R_n=r_n(y_n),\qquad A_n=r_n(p_n^{\rm in}),\qquad X_n=-A_n,
 \qquad c_1=\frac{1-q}{2},
 \label{eq:terminal-modal-chronology-residuals}
\end{equation}
and, on coordinates shared by two consecutive prefixes,
\begin{equation}
 D_n(j)=X_n(j)-c_1X_{n-1}(j),\qquad
 2\le n\le m,\quad 0\le j\le n-2.
 \label{eq:terminal-modal-shared-sign}
\end{equation}

\begin{lemma}[Exact chronology reduction; proved here]
\label{lem:terminal-modal-chronology-reduction}
For the named parameters, if
$D_n(j)\ge0$ throughout the range in
\eqref{eq:terminal-modal-shared-sign}, then the projected execution agrees
with the linear candidate at every proper prefix and satisfies
\cref{lem:terminal-modal-admission-chronology}.  Thus the only remaining
proper-prefix sign target is the uniform shared-coordinate inequality
\eqref{eq:terminal-modal-shared-sign}; the raw-frontier, post-residual, and
outside-gate implications require no additional asymptotic hypothesis.
\end{lemma}

\begin{proof}
Write $\sigma_n=\widetilde p_n(n-1)$ for the raw frontier value.  Suppose the
candidate has agreed with the projected execution through prefix $n-1$ and
$R_{n-1}\le0$.  At entry to prefix $n\ge2$, exact center transport gives
\[
 y_n(n-1)=\frac{qP_n}{1+q},\qquad
 y_n(n-2)=\frac{2\sigma_{n-1}
                 +q(b_n-b_{n-1})f_{n-2}}{1+q}.
\]
Since $\sigma_{n-1}\ge y_{n-1}(n-2)=qP_{n-1}/(1+q)$,
\cref{lem:terminal-modal-frontier-bounds} and
$f_{n-2}/f_{n-1}\ge\chi$ imply
\begin{equation}
 \frac{y_n(n-2)}{y_n(n-1)}
 \ge\frac2{1+q}+\chi=\frac{3-q}{1+q}.
 \label{eq:terminal-modal-frontier-ratio}
\end{equation}
The degree-two frontier row of $B_n=I-\widehat Q_n$ therefore gives
\begin{align}
 \frac{\sigma_n}{q^3}
 &\ge \frac{(1-q)(5+q)}{4(1+q)}\frac{P_n}{q^2}-\frac15
 \notag\\
 &\ge \frac{1215}{1088}\frac{3488}{1921}-\frac15
 =:L=\frac{596861}{326570}.
 \label{eq:terminal-modal-raw-frontier-lower}
\end{align}
Here $(1-q)(5+q)/(4(1+q))$ is decreasing on this parameter interval, so its
minimum is the displayed $1215/1088$.
The base value $\sigma_1=q^2-q^3/5$ satisfies the same lower bound.  Since
\[
 L-\frac{2048}{1275}
 =\frac{1084499}{4898550}>0,
 \qquad
 \frac4{5\eta}\le\frac{2048}{1275},
\]
the next degree-two outside residual
$q^3/5-\eta\sigma_n/2$ is strictly below $-q^3/5$.  At $n=m$ the next row is
the degree-one endpoint, whose threshold is only $2/(5\eta)$, so the same
bound is more than sufficient.

It remains to prove that $R_n\le0$.  On every shared coordinate, the exact
transport identity and the fact that the optimum displacement is harmonic on
old rows give
\begin{equation}
 R_n(j)=\frac{2A_n(j)-(1-q)A_{n-1}(j)}{1+q}
       =-\frac{2D_n(j)}{1+q}\le0.
 \label{eq:terminal-modal-shared-average-residual}
\end{equation}
At the new frontier, direct evaluation of the padded old error and the
optimum displacement instead gives
\begin{equation}
 R_n(n-1)=\frac{-\eta\sigma_{n-1}
                 +q\eta P_{n-1}/2+q^3/5}{1+q}.
 \label{eq:terminal-modal-frontier-average-residual}
\end{equation}
Using \eqref{eq:terminal-modal-raw-frontier-lower},
$P_{n-1}/q^2\le33/16$, and $\eta\ge255/512$ shows that the numerator divided
by $q^3$ is at most
\[
 \frac{255}{512}\left(-L+\frac{33}{32}\right)+\frac15
 =-\frac{30946901}{157368320}<0.
\]
The base residual is $R_1=-q^2+q^3/5<0$, completing the induction for the
average residual.

Finally, $B_n$ is entrywise nonnegative and the affine residual identity is
\[
 r_n(\widetilde p_n)=B_nR_n\le0.
\]
Hence the safe correction is exactly zero.  The invariant
$v_n^{\rm in}\ge p_n^{\rm in}$ starts at equality and propagates: from
$R_n\le0$ one has
$\widetilde p_n\ge y_n\ge p_n^{\rm in}$, the center update is at least
$\widetilde p_n$, and the optimum displacement is strictly positive.
Thus $y_n$ and the raw point are strictly positive for $n\ge2$; the base raw
point is positive directly.  Projection is inactive, the strict outside
test admits the unique adjacent singleton, and the projected execution
continues to agree with the candidate.
\end{proof}

The remaining sign does not follow from a bare positivity-preserving
recurrence.  Indeed, on the even extension and in the linear regime,
\eqref{eq:terminal-modal-directional-factorization} gives exactly, with the
boundary defect $F_{n+1}$ defined in
\eqref{eq:terminal-modal-boundary-source} below,
\begin{equation}
 D_{n+1}^{\rm ev}
 =c_1(I+\mathcal S_++\mathcal S_-)D_n^{\rm ev}
  -c_1^2X_{n-1}^{\rm ev}-F_{n+1}^{\rm ev}.
 \label{eq:terminal-modal-shared-sign-recurrence}
\end{equation}
Both the global negative term and the positive $U_n,2U_n$ entries of
$F_{n+1}^{\rm ev}$ obstruct a direct maximum-principle induction.  A proof of
\eqref{eq:terminal-modal-shared-sign} therefore still needs a quantitative
spatial cone or the signed moving-frontier cancellation; the finite screen is
not such a proof.

For the remainder of this section assume the conclusion of
\cref{lem:terminal-modal-admission-chronology}.

Under that hypothesis, compute the changing-face defect exactly.  Write $p_n^{\rm in}$ for
the literal state immediately after prefix $n$ is admitted, and define
\begin{equation}
 e_n=p_n^{\rm in}-p_n^\star,\qquad
 A_n=\widehat Q_ne_n=r_n(p_n^{\rm in}),\qquad
 \Delta_n=b_n-b_{n-1}.
 \label{eq:terminal-modal-entry-error}
\end{equation}
Vectors on shorter prefixes are zero-extended whenever they occur in a
longer-prefix formula.  If $B_n=I-\widehat Q_n$, projection inactivity, zero
correction, and exact center transport give
\begin{equation}
 e_{n+1}
 =\left[\frac{B_n}{1+q}
 \left(2e_n-(1-q)e_{n-1}+\Delta_nf\right);0\right]
 -\Delta_{n+1}f,
 \qquad 2\le n\le m-1.
 \label{eq:terminal-modal-changing-face-error}
\end{equation}
Indeed, if $e_n^v=v_n^{\rm in}-p_n^\star$, transport gives
$e_n^v=e_n+(1-q)(e_n-e_{n-1})/q+\Delta_nf/q$; substituting this in the
accelerated average and then zero-padding the new primal coordinate proves
\eqref{eq:terminal-modal-changing-face-error}.

Zero-extend each $A_n$ on the nonnegative half-line and define
\begin{equation}
 F_{n+1}:=E_{n+1}-\mathcal R_-E_n.
 \label{eq:terminal-modal-boundary-source}
\end{equation}

\begin{lemma}[Exact changing-face boundary source under the open chronology]
\label{lem:terminal-modal-boundary-source}
Put
\begin{equation}
 c_1=\frac{1-q}{2},\qquad c_2=\frac{(1-q)^2}{4},\qquad
 P_n=p_n^\star(n-1),
 \label{eq:terminal-modal-source-constants}
\end{equation}
and, for $4\le n\le m-1$, define
\begin{align}
 U_n&=\frac{q\eta^2}{4(1+q)}
 \left(P_{n-1}-\frac{2q\rho}{\eta}\right),
 \label{eq:terminal-modal-source-U}\\
 W_n&=U_n+c_1P_n-
 \frac{a\Delta_n}{1+q}f_n
 +\frac\eta2\left(\frac{\Delta_n}{1+q}-\Delta_{n+1}\right)f_{n+1}.
 \label{eq:terminal-modal-source-W}
\end{align}
Then
\begin{equation}
 \operatorname{supp}F_{n+1}\subseteq\{0,n-2,n-1,n\},
 \label{eq:terminal-modal-source-support}
\end{equation}
and its possibly nonzero entries are exactly
\begin{align}
 F_{n+1}(0)&=c_1A_n(1)-c_2A_{n-1}(1),
 \label{eq:terminal-modal-source-seed}\\
 F_{n+1}(n-2)&=U_n,
 &F_{n+1}(n-1)&=2U_n,
 &F_{n+1}(n)&=W_n.
 \label{eq:terminal-modal-source-frontier}
\end{align}
The cases $n<4$ have overlapping seed and frontier stencils and are obtained
by direct substitution in
\eqref{eq:terminal-modal-changing-face-error}.

At the final transition $n=m$, the support statement remains valid and the
first three displayed entries are unchanged.  Only the last entry changes.
Writing $A_m^{\rm fr}=A_m(m-1)$, it is
\begin{equation}
 \begin{split}
 W_m^{\rm end}={}&c_1A_m^{\rm fr}
 +\frac{(1-q)^2}{2}P_m
 -\frac{\eta(1-q)^2}{4}P_{m-1}\\
 &-\frac{\Delta_m\eta^2}{1+q}
 \left(f_{m-1}+\frac12f_{m-2}\right)+\frac{q^3}{5}.
 \end{split}
 \label{eq:terminal-modal-source-final-endpoint}
\end{equation}
This is the sole extra term caused by changing the new vertex's ambient
degree from two to one.
\end{lemma}

\begin{proof}
On rows $j\ge1$, let
\[
 H_0=aI-\frac\eta2(\mathcal S_++\mathcal S_-),
 \qquad
 B_0=\eta I+\frac\eta2(\mathcal S_++\mathcal S_-).
\]
The actual seed row differs only in that its right-neighbor coefficient is
doubled.  Substitute
\eqref{eq:terminal-modal-changing-face-error} into
$A_{n+1}=\widehat Q_{n+1}e_{n+1}$ and subtract the two fixed-face stencil
terms in \eqref{eq:terminal-modal-boundary-source}.  The $e_n$ and
$e_{n-1}$ coefficients cancel on every row outside the displayed stencils.
At the seed, the remaining expression is
\[
 -\frac{\eta}{2(1+q)}
  (\eta e_n(0)-2ae_n(1)+\eta e_n(2))
 +\frac{(1-q)\eta}{4(1+q)}
  (\eta e_{n-1}(0)-2ae_{n-1}(1)+\eta e_{n-1}(2)).
\]
The two parenthesized quantities are $-2A_n(1)$ and
$-2A_{n-1}(1)$, respectively, which proves
\eqref{eq:terminal-modal-source-seed}.

At entry, the newest primal coordinate is zero, so
$e_n(n-1)=-P_n$ and $e_{n-1}(n-2)=-P_{n-1}$.  Direct evaluation of the last
three rows first gives
\[
 \begin{split}
 \widetilde U_n={}&
 -\frac{(1-q)\eta^2}{4(1+q)}P_{n-1}
 +\frac{\Delta_n\eta^2}{4(1+q)}f_n,\\
 F_{n+1}(n-2)={}&\widetilde U_n,\qquad
 F_{n+1}(n-1)=2\widetilde U_n,\\
 F_{n+1}(n)={}&\widetilde U_n+c_1P_n
 -\frac{a\Delta_n}{1+q}f_n
 +\frac\eta2\left(\frac{\Delta_n}{1+q}-\Delta_{n+1}\right)f_{n+1}.
 \end{split}
\]
The old optimum has ghost value zero one coordinate before the new
frontier.  Applying its affine recurrence at the next ghost coordinate gives
\begin{equation}
 \Delta_nf_n=P_{n-1}-\frac{2q^2\rho}{\eta}.
 \label{eq:terminal-modal-source-ghost-identity}
\end{equation}
Consequently $\widetilde U_n=U_n$, proving
\eqref{eq:terminal-modal-source-frontier}.

On the last transition the old face still has the same degree-two frontier,
whereas the new endpoint row is $az_m-\eta z_{m-1}$.  Repeating only this
last-row calculation and using
$\widehat Q_{m+1}(p_{m+1}^\star-[p_m^\star;0])
=(\eta P_m-q^3/5)e_m$ yields
\eqref{eq:terminal-modal-source-final-endpoint}; all earlier rows are
unchanged.
\end{proof}

The seed source is an artifact of zero extension.  It disappears under the
even extension
\begin{equation}
 A_n^{\rm ev}(j)=
 \begin{cases}
  A_n(|j|),& |j|<n,\\
  0,& |j|\ge n.
 \end{cases}
 \label{eq:terminal-modal-even-extension}
\end{equation}
Indeed, the reflected neighbor supplies precisely the two terms in
\eqref{eq:terminal-modal-source-seed}.  Thus the defect of
$A_n^{\rm ev}$ on $\mathbb Z$ consists only of the three positive-frontier
entries in \eqref{eq:terminal-modal-source-frontier} and their reflections.

This representation exposes a second cancellation.  Put
\begin{equation}
 s_n=\frac{f_{n-1}}{f_n},\qquad
 D=\frac{2q^2\rho}{\eta},\qquad M_n=W_n+3U_n.
 \label{eq:terminal-modal-source-mass-parameters}
\end{equation}

\begin{lemma}[Factored frontier transform under the open chronology]
\label{lem:terminal-modal-source-transform}
For $4\le n\le m-1$ and $z=e^{i\theta}$, the Fourier transform of the even
source is
\begin{equation}
 \widehat F_{n+1}^{\rm ev}(\theta)
 =2\Re\left[z^{n-2}
 \left\{U_n(1-z)(1+3z)+M_nz^2\right\}\right].
 \label{eq:terminal-modal-source-transform}
\end{equation}
Moreover
\begin{equation}
 M_n=\frac{q(1-q)}4(2\eta-s_n)P_{n-1}
 +\frac q4\bigl((1-q)s_n+2\eta\bigr)D.
 \label{eq:terminal-modal-source-mass}
\end{equation}
For the named parameters, $U_n$ is nonincreasing in $n$ and, uniformly in
these $n$,
\begin{equation}
 0<U_n\le\frac{3}{40}q^3,\qquad
 |M_n|\le\frac32q^4.
 \label{eq:terminal-modal-source-bounds}
\end{equation}
\end{lemma}

\begin{proof}
The three positive-frontier entries have polynomial
$z^{n-2}(U_n+2U_nz+W_nz^2)$.  Reflection takes twice its real part, and
$1+2z-3z^2=(1-z)(1+3z)$, proving
\eqref{eq:terminal-modal-source-transform}.

The identities
\[
 \Delta_nf_n=P_{n-1}-D,\qquad
 P_n=\Delta_nf_{n-1}=s_n(P_{n-1}-D),\qquad
 s_n+\frac{f_{n+1}}{f_n}=\frac{2a}{\eta}
\]
reduce $W_n+3U_n$ exactly to
\eqref{eq:terminal-modal-source-mass}.

The monotonicity and positivity needed here follow from
\cref{lem:terminal-modal-frontier-bounds}.  Evaluating
\eqref{eq:terminal-modal-frontier-optimum} at $t=1$ and using
\eqref{eq:terminal-modal-source-U} give
\[
 \frac{U_n}{q^3}
 \le\frac{\eta^2}{4(1+q)}
 \left(\frac{6/5}{1-q}+\frac{4/5}{1+q}-\frac{2}{5\eta}\right)
 =\frac{(1-q)(6/5+2q/5)}{16}\le\frac3{40}.
\]
Finally $\chi<s_n\le1$,
$|2\eta-s_n|\le2q$, $P_{n-1}\le33q^2/16$, and
$D\le13q^3/16$ for $q\le1/16$.  Substitution in
\eqref{eq:terminal-modal-source-mass} gives
$|M_n|\le(33/32+13/32)q^4<3q^4/2$.
\end{proof}

Let $\mathcal A_n(\theta)$ be the Fourier transform of
$A_n^{\rm ev}$.  With $\phi=\theta/2$ and
$r=(1-q)\cos\phi$, the exact scalar recurrence is
\begin{equation}
 \mathcal A_{n+1}=2r\cos\phi\,\mathcal A_n-r^2\mathcal A_{n-1}
 +\widehat F_{n+1}^{\rm ev}.
 \label{eq:terminal-modal-driven-transform}
\end{equation}
Consequently a source injected at time $t$ reaches time $N$ with multiplier
\begin{equation}
 r^{N-t}\frac{\sin((N-t+1)\phi)}{\sin\phi}.
 \label{eq:terminal-modal-source-green}
\end{equation}
The factor $1-z$ in
\eqref{eq:terminal-modal-source-transform} cancels the apparent small-sine
loss in this Green kernel.  What remains open is a uniform signed
summation-by-parts estimate for the slowly varying trace $U_n$, together with
the $M_n$ and final-endpoint contributions, sharp enough to prove
\eqref{eq:terminal-modal-entry-profiles}.  Absolute-value summation discards
the moving-frontier phase and does not deliver the stated constants.

For reference, this signed estimate can be stated as two exact scalar
targets.  At an even full-path mode put
$\theta=2s\pi/m$, let
\begin{equation}
 Y_{m+1}(\theta)
 =\mathcal A_{m+1}(\theta)-A_{m+1}(m),
 \qquad
 \kappa_m=\eta P_m-\frac{q^3}{5}.
 \label{eq:terminal-modal-source-scalar-data}
\end{equation}
The endpoint subtraction in $Y_{m+1}$ changes the reflected weight two to
the full path's ambient endpoint weight one.

\begin{lemma}[Exact scalar position/velocity reduction under the open chronology]
\label{lem:terminal-modal-source-scalar-reduction}
For every even internal mode $h=2s$,
\begin{align}
 mC_{2s}&=Y_{m+1}(\theta),
 \label{eq:terminal-modal-source-position-scalar}\\
 mV_{2s}&=\frac1q\left[
 Y_{m+1}(\theta)+\kappa_m
 -(1-q)\bigl(\mathcal A_m(\theta)-\eta e_m(m-1)\bigr)
 \right].
 \label{eq:terminal-modal-source-velocity-scalar}
\end{align}
Hence the two open profile inequalities are equivalent to
\begin{align}
 \left|Y_{m+1}(\theta)-mG_{2s}\right|&\le\frac{q^2}{256},
 \label{eq:terminal-modal-source-position-target}\\
 \left|Y_{m+1}(\theta)+\kappa_m
 -(1-q)\bigl(\mathcal A_m(\theta)-\eta e_m(m-1)\bigr)
 +\frac q5mG_{2s}\right|&\le\frac{q^3}{4}.
 \label{eq:terminal-modal-source-velocity-target}
\end{align}
\end{lemma}

\begin{proof}
The first identity is the definition of the degree-weighted cosine
coefficient, using even extension and correcting the far-endpoint weight.
For the second, let $e_m^{\rm out}$ be the centered error after the last
proper-prefix step.  At full-face entry,
\[
 e_{m+1}=[e_m^{\rm out};0]
 -(p_{m+1}^\star-[p_m^\star;0]),
 \qquad
 e_{m+1}^v=[e_m^{\rm out};0]
 +\frac{1-q}{q}([e_m^{\rm out};0]-e_m).
\]
The full operator applied to the final optimum displacement is
$\kappa_m$ at the far endpoint and zero elsewhere.  Its application to the
zero extension of $e_m$ equals $A_m$ on the old prefix and
$-\eta e_m(m-1)$ at the far endpoint.  Taking the degree-weighted even-mode
coefficient gives
\eqref{eq:terminal-modal-source-velocity-scalar}.  Multiplying the profile
inequalities by $q^2$ and then, in the velocity case, by $q$ proves the two
equivalences.
\end{proof}

For $U_0=3q^3/40$, split the scalar target into: (i) the homogeneous and
$n<4$ response minus the ideal packet; (ii) constant-$U_0$ response with
$M_n=0$; (iii) $U_n-U_0$ response; (iv) $M_n$ response; and (v) final
degree/endpoint correction.  Assign $qmG_{2s}/5$ to the first velocity piece
and $\kappa_m+(1-q)\eta e_m(m-1)$ to the fifth.  Linearity makes the five sums
exactly \eqref{eq:terminal-modal-source-position-target} and the expression
inside \eqref{eq:terminal-modal-source-velocity-target}, respectively.

The open task is a uniform signed bound on these sums.  At velocity scale the
constant-$U$ term has a nonzero order-$q^3$ limit that cancels against the
base and endpoint terms, so a triangle inequality cannot suffice.  One needs
that leading response, summation by parts for $U_n-U_0$, and bounds on the
$M_n=O(q^4)$ and endpoint remainders.
